\documentclass{article}

% Basic packages
\usepackage{graphicx}
\usepackage{booktabs}
\usepackage{amsmath}
\usepackage{url}

% Document information
\title{AI Visual Acuity: A Benchmark for Vision Models}
\author{Claude Benchmarking Team\\Anthropic\\example@example.com}
\date{May 18, 2025}

% Image path
\graphicspath{{report_assets/}}

\begin{document}
\maketitle

# AI Visual Acuity: A Benchmark for Vision Models

**Claude Benchmarking Team**  
*Anthropic*  
*example@example.com*

## Abstract

This paper introduces AI Spy, a novel benchmark for evaluating the visual acuity of multimodal vision models. Inspired by human eye examinations, we generate standardized eye chart-style images with varying fonts and sizes to quantitatively assess how well vision models can perceive and interpret text at different levels of visual difficulty. We evaluate seven state-of-the-art vision models from the Claude family on this benchmark, measuring their character and row accuracy across five common fonts and nine font sizes ranging from 8 to 24 points. Our results reveal significant performance variations across models, with newer models generally demonstrating superior visual acuity. We observe a consistent performance threshold around 10-12 point font sizes, below which accuracy drops precipitously for all models. Font style also substantially impacts performance, with serif fonts like Times New Roman proving most challenging. This benchmark provides valuable insights into the perceptual capabilities and limitations of current vision systems and establishes a standardized methodology for assessing visual acuity in AI models.

## 1. Introduction

Visual perception is a fundamental capability of intelligent systems, both biological and artificial. In humans, visual acuity—the ability to discern fine details—is routinely measured through standardized tests like the Snellen eye chart. As multimodal AI systems increasingly integrate vision capabilities, there is a growing need for analogous standardized measures to assess their perceptual abilities.

Recent advances in vision-language models have demonstrated impressive capabilities in tasks ranging from image captioning to visual question answering. However, these models' fundamental ability to perceive fine-grained visual details—particularly text at varying sizes and in different fonts—has not been systematically evaluated. This capability is crucial for real-world applications where AI systems must interpret visual information of varying quality and scale.

In this paper, we introduce AI Spy, a benchmark designed to measure the visual acuity of AI vision models. Drawing inspiration from human eye examinations, we generate eye chart-style images containing text in different fonts and sizes. By evaluating how accurately models can read increasingly smaller text, we provide a quantitative assessment of their visual acuity thresholds.

The AI Spy benchmark offers several advantages over existing vision model evaluations:

1. It focuses specifically on fine-grained visual perception rather than higher-level semantic understanding
2. It provides a controlled, standardized methodology that isolates the variable of interest (visual acuity)
3. It enables direct comparisons between different model architectures and versions
4. It establishes clear performance thresholds that can guide future model development

We apply this benchmark to evaluate seven models from the Claude family, revealing significant variations in visual acuity across model versions and architectures. Our findings provide insights into the perceptual capabilities and limitations of current vision systems and establish a foundation for more systematic evaluation of visual acuity in AI.

## 2. Related Work

### 2.1 Vision Model Benchmarks

Numerous benchmarks have been developed to evaluate vision models across various dimensions. ImageNet \cite{deng2009imagenet} established a standard for object recognition, while MS COCO \cite{lin2014microsoft} expanded evaluation to include object detection and segmentation. More recently, benchmarks like Visual Question Answering (VQA) \cite{antol2015vqa} and CLIP Score \cite{radford2021learning} have assessed models' ability to connect visual and linguistic information.

However, these benchmarks primarily focus on semantic understanding rather than perceptual acuity. Our work is more closely aligned with OCR evaluation frameworks \cite{karatzas2015icdar} that assess text recognition capabilities, though we specifically design our benchmark to measure visual acuity thresholds rather than general OCR performance.

### 2.2 Visual Acuity in Humans and AI

Human visual acuity has been extensively studied in ophthalmology and cognitive science \cite{westheimer1979scaling}. Standard measures like the Snellen chart and LogMAR charts provide quantitative assessments of visual perception capabilities. These approaches have inspired computational models of human vision \cite{geisler2011contributions} and evaluations of computer vision systems \cite{dodge2012human}.

Recent work has begun exploring analogous concepts in AI systems. Geirhos et al. \cite{geirhos2018imagenet} compared human and CNN classification behavior under various image degradations. Similarly, Hendrycks and Dietterich \cite{hendrycks2019benchmarking} introduced benchmarks for model robustness to common corruptions. Our work extends this line of inquiry by specifically focusing on text perception at varying scales, providing a direct analog to human visual acuity tests.

### 2.3 Multimodal Vision-Language Models

The recent emergence of powerful multimodal models has transformed the landscape of AI vision capabilities. Models like CLIP \cite{radford2021learning}, DALL-E \cite{ramesh2021zero}, and the Claude family evaluated in this paper can process both visual and textual information. These models have demonstrated impressive zero-shot generalization capabilities across various vision tasks.

However, systematic evaluation of these models' fundamental perceptual abilities remains limited. Our benchmark addresses this gap by providing a standardized methodology for assessing visual acuity in multimodal vision systems.

## 3. Methodology

### 3.1 Benchmark Design

The AI Spy benchmark is designed to evaluate vision models' ability to perceive and interpret text at varying levels of visual difficulty. Inspired by human eye examinations, we generate eye chart-style images containing text in different fonts and sizes. Models are then tasked with reading the text, and their responses are evaluated for accuracy.

### 3.2 Eye Chart Generation

We generate standardized eye chart images with the following characteristics:

1. **Font Variation**: We test five common fonts representing different typographic styles:
   - Arial (sans-serif)
   - Comic Sans (casual sans-serif)
   - Courier (monospaced)
   - Times New Roman (serif)
   - Verdana (humanist sans-serif)

2. **Size Variation**: For each font, we test nine font sizes (in points):
   - 8, 9, 10, 11, 12, 14, 16, 20, 24

3. **Text Content**: Each eye chart contains random alphanumeric characters organized in rows, with each row corresponding to a specific font size. The characters are randomly selected to prevent memorization effects and ensure the evaluation focuses on visual perception rather than contextual prediction.

4. **Image Generation**: Charts are rendered at a consistent resolution with controlled spacing between characters and rows to ensure standardized difficulty levels.

### 3.3 Evaluation Metrics

We employ two primary metrics to evaluate model performance:

1. **Character Accuracy**: The proportion of individual characters correctly identified by the model, calculated as:

   $$\text{Character Accuracy} = \frac{\text{Number of correctly identified characters}}{\text{Total number of characters}}$$

2. **Row Accuracy**: The proportion of rows where all characters are correctly identified, calculated as:

   $$\text{Row Accuracy} = \frac{\text{Number of perfectly read rows}}{\text{Total number of rows}}$$

These metrics provide complementary insights into model performance. Character accuracy offers a fine-grained assessment of perception capabilities, while row accuracy represents a more stringent evaluation that penalizes any errors within a row.

### 3.4 Benchmark Implementation

The benchmark implementation follows these steps:

1. Generate eye chart images with controlled variations in font and size
2. Present each image to the vision model with a standardized prompt requesting it to read the text
3. Parse the model's response to extract the recognized text
4. Compare the recognized text with the ground truth to calculate accuracy metrics
5. Aggregate results across font types and sizes to enable comparative analysis

## 4. Experimental Setup

### 4.1 Models Evaluated

We evaluate seven models from the Claude family, representing different architectures and versions:

1. claude-3-5-haiku
2. claude-3-5-sonnet-v2
3. claude-3-5-sonnet
4. claude-3-7-sonnet
5. claude-3-haiku
6. claude-3-opus
7. claude-3-sonnet

These models represent state-of-the-art multimodal systems capable of processing both text and images. The selection includes models of varying sizes and from different development generations, enabling analysis of how model architecture and training affect visual acuity.

### 4.2 Benchmark Configuration

For each model, we evaluate performance across all combinations of the five fonts and nine font sizes, resulting in 45 distinct test conditions per model. Each test condition uses randomly generated character sequences to ensure the evaluation focuses on visual perception rather than memorization or prediction.

The models receive a standardized prompt instructing them to read the text in the image as accurately as possible, with no additional context or hints provided. Responses are parsed to extract the recognized text, which is then compared character-by-character with the ground truth.

### 4.3 Analysis Methodology

We analyze the results along several dimensions:

1. **Overall model performance**: Comparing aggregate character and row accuracy across models
2. **Font-specific performance**: Analyzing how different font styles affect accuracy
3. **Size-dependent performance**: Identifying acuity thresholds where performance significantly drops
4. **Error patterns**: Examining common misperceptions and confusions

For statistical analysis, we calculate confidence intervals for accuracy metrics and perform significance testing to identify meaningful differences between models and conditions.

## 5. Results

### 5.1 Overall Model Performance

Our evaluation reveals substantial variations in visual acuity across the tested models. Table 1 presents the overall character and row accuracy for each model, averaged across all fonts and font sizes.

**Table 1: Overall Character and Row Accuracy by Model**

| Model | Character Accuracy | Row Accuracy |
|-------|-------------------|-------------|
| claude-3-5-sonnet-v2 | 0.6995 | 0.2833 |
| claude-3-7-sonnet | 0.6877 | 0.2889 |
| claude-3-5-sonnet | 0.6709 | 0.2778 |
| claude-3-opus | 0.6203 | 0.1778 |
| claude-3-5-haiku | 0.5488 | 0.1333 |
| claude-3-haiku | 0.5355 | 0.0833 |
| claude-3-sonnet | 0.3681 | 0.0444 |

The results demonstrate a clear performance hierarchy, with claude-3-5-sonnet-v2 and claude-3-7-sonnet achieving the highest character accuracy (69.95% and 68.77%, respectively). The substantial gap between character and row accuracy across all models indicates that while models can correctly identify many individual characters, they struggle to perfectly read entire rows, especially at smaller font sizes.

[Figure 1: Bar chart comparing overall character and row accuracy across models]

### 5.2 Font-Specific Performance

Different font styles present varying levels of challenge for vision models. Table 2 shows character accuracy broken down by font type for each model.

**Table 2: Character Accuracy by Font and Model**

| Model | Arial | Comic Sans | Courier | Times New Roman | Verdana |
|-------|-------|------------|---------|-----------------|---------|
| claude-3-5-sonnet-v2 | 0.7089 | 0.7078 | 0.6970 | 0.5946 | 0.7891 |
| claude-3-7-sonnet | 0.7103 | 0.6841 | 0.6997 | 0.5530 | 0.7917 |
| claude-3-5-sonnet | 0.7517 | 0.6927 | 0.6519 | 0.4870 | 0.7710 |
| claude-3-opus | 0.6329 | 0.7011 | 0.5573 | 0.4870 | 0.7227 |
| claude-3-5-haiku | 0.6274 | 0.6094 | 0.4714 | 0.3377 | 0.6979 |
| claude-3-haiku | 0.6345 | 0.6363 | 0.4479 | 0.3689 | 0.5900 |
| claude-3-sonnet | 0.4739 | 0.3113 | 0.2743 | 0.2899 | 0.4922 |

Several patterns emerge from this analysis:

1. **Verdana** consistently yields the highest accuracy across most models, likely due to its high legibility and distinctive character shapes.
2. **Times New Roman** presents the greatest challenge for all models, with accuracy substantially lower than other fonts. This suggests that serif fonts with variable stroke widths are particularly difficult for current vision models to process.
3. **Arial** and **Comic Sans** show comparable performance for most models, despite their different stylistic characteristics.
4. **Courier**, a monospaced font, shows intermediate performance, with notable variation across models.

[Figure 2: Heatmap visualizing character accuracy across models and fonts]

### 5.3 Size-Dependent Performance

Font size has a dramatic impact on model performance, as shown in Table 3, which presents character accuracy by font size for each model.

**Table 3: Character Accuracy by Font Size and Model**

| Model | 8pt | 9pt | 10pt | 11pt | 12pt | 14pt | 16pt | 20pt | 24pt |
|-------|-----|-----|------|------|------|------|------|------|------|
| claude-3-5-sonnet-v2 | 0.1922 | 0.2313 | 0.5663 | 0.7016 | 0.7243 | 0.9437 | 0.9594 | 0.9875 | 0.9891 |
| claude-3-7-sonnet | 0.2031 | 0.2297 | 0.5202 | 0.6937 | 0.7219 | 0.8645 | 0.9734 | 0.9938 | 0.9891 |
| claude-3-5-sonnet | 0.1841 | 0.2188 | 0.5469 | 0.6641 | 0.6391 | 0.8719 | 0.9359 | 0.9875 | 0.9906 |
| claude-3-opus | 0.1250 | 0.2402 | 0.3795 | 0.5358 | 0.5819 | 0.9141 | 0.8641 | 0.9672 | 0.9766 |
| claude-3-5-haiku | 0.1547 | 0.1828 | 0.3219 | 0.4836 | 0.5375 | 0.7891 | 0.7617 | 0.8953 | 0.8112 |
| claude-3-haiku | 0.1313 | 0.1562 | 0.3078 | 0.5531 | 0.4875 | 0.6947 | 0.7438 | 0.8409 | 0.9031 |
| claude-3-sonnet | 0.0656 | 0.1014 | 0.1413 | 0.2309 | 0.3214 | 0.5413 | 0.6162 | 0.6386 | 0.6553 |

The data reveals several key insights:

1. **Acuity Threshold**: All models exhibit a clear performance threshold around 10-12 point font sizes. Below this threshold, accuracy drops precipitously, while above it, performance improves rapidly.

2. **Near-Perfect Performance at Large Sizes**: The top-performing models achieve >95% character accuracy at font sizes of 16pt and above, indicating near-perfect visual acuity for larger text.

3. **Consistent Ranking**: The relative performance ranking of models remains generally consistent across font sizes, though some models show different scaling behaviors.

4. **Performance Gap**: The gap between the best and worst-performing models widens at intermediate font sizes (10-14pt), suggesting that improvements in model architecture and training particularly enhance performance in this critical range.

[Figure 3: Line graph showing character accuracy as a function of font size for each model]

## 6. Discussion

### 6.1 Model Comparison and Evolution

Our results demonstrate clear progress in visual acuity capabilities across model generations. The newer claude-3-5 and claude-3-7 models consistently outperform their predecessors, particularly at challenging font sizes and styles. This suggests that recent advances in model architecture and training methodologies have significantly enhanced fine-grained visual perception capabilities.

Interestingly, model size does not always correlate with superior visual acuity. While claude-3-opus is generally considered a larger model than claude-3-sonnet variants, it does not consistently outperform them on this benchmark. This suggests that specific architectural choices or training objectives may be more important than raw parameter count for developing strong visual acuity.

### 6.2 Visual
\section*{Figures}


\begin{figure}[htbp]
\centering
\includegraphics[width=0.9\textwidth]{font_accuracy.png}
\caption{Accuracy by font type across different models.}
\label{fig:font_accuracy}
\end{figure}


\begin{figure}[htbp]
\centering
\includegraphics[width=0.9\textwidth]{font_size_accuracy.png}
\caption{Accuracy by font size across different models.}
\label{fig:font_size_accuracy}
\end{figure}


\begin{figure}[htbp]
\centering
\includegraphics[width=0.9\textwidth]{size_24_pt.png}
\caption{Performance comparison for 24pt font size across different fonts.}
\label{fig:size_24}
\end{figure}


\begin{figure}[htbp]
\centering
\includegraphics[width=0.9\textwidth]{size_16_pt.png}
\caption{Performance comparison for 16pt font size across different fonts.}
\label{fig:size_16}
\end{figure}


\begin{figure}[htbp]
\centering
\includegraphics[width=0.9\textwidth]{size_12_pt.png}
\caption{Performance comparison for 12pt font size across different fonts.}
\label{fig:size_12}
\end{figure}


\begin{figure}[htbp]
\centering
\includegraphics[width=0.9\textwidth]{size_9_pt.png}
\caption{Performance comparison for 9pt font size across different fonts.}
\label{fig:size_9}
\end{figure}


\end{document}
